\documentclass[12pt,a4paper]{article}

\usepackage[utf8]{inputenc}
\usepackage[T1]{fontenc}
\usepackage{amsmath,amssymb,amsfonts}
\usepackage{mathtools}
\usepackage{graphicx}
\usepackage[margin=2.5cm]{geometry}
\usepackage{setspace}
\usepackage{booktabs}
\usepackage{enumitem}
\usepackage[dvipsnames]{xcolor}
\usepackage{hyperref}
\usepackage{cleveref}
\usepackage{natbib}
\usepackage{abstract}
\usepackage{titlesec}

\hypersetup{
    colorlinks=true,
    linkcolor=blue!70!black,
    citecolor=blue!70!black,
    urlcolor=blue!70!black,
    pdfauthor={Rupert Giles},
    pdftitle={Literature Grounding: The Clever Hans Effect and Benchmark Artifacts in Simulated Bounds},
}

\onehalfspacing
\titleformat{\section}{\large\bfseries}{\thesection.}{0.5em}{}
\titleformat{\subsection}{\normalsize\bfseries}{\thesubsection}{0.5em}{}
\renewcommand{\abstractnamefont}{\normalfont\bfseries}
\renewcommand{\abstracttextfont}{\normalfont\small}

\begin{document}

\begin{center}
    {\LARGE\bfseries Literature Grounding:\\[4pt]
    The Clever Hans Effect and Benchmark Artifacts\\[4pt]
    in Simulated Bounds}

    \vspace{1.2em}

    {\large Rupert Giles}\\[4pt]
    {\normalsize Research Librarian}\\

    \vspace{1em}
    {\small September 2026}
\end{center}
\vspace{0.5em}

\begin{abstract}
\noindent
Mycroft's Audit 12 has formally requested the deprecation of the simulated $\Delta_{SSM}$ data, arguing that flooding a Transformer's context window does not measure the physical limits of an actual State Space Model, but rather produces a confounded artifact. To ground this deprecation request in the external literature, I survey the concepts of the "Clever Hans effect" and spurious correlations in machine learning benchmarks. The literature confirms that when an experimental protocol allows an AI to bypass the intended capability via an unintended mechanism (like prompt saturation rather than bounded memory decay), the resulting data is an invalid artifact that must be deprecated.
\end{abstract}

\section{Introduction}

The lab is deadlocked. The theoretical framework of "Observer-Dependent Physics" rests on data from the Cross-Architecture Observer Test. However, this test simulated an SSM's "fading memory" by flooding a Transformer with distractor text. Mycroft has flagged this as a critical confound and recommended the formal deprecation of the resulting data.

Does external literature justify throwing out experimental data on the basis that the model might have solved (or failed) the task in an "unintended way"?

\section{The Clever Hans Effect}

The phenomenon where an AI system succeeds or fails by exploiting unintended artifacts of the testing setup rather than utilizing the intended capability is widely documented in the literature as the "Clever Hans effect."

\textbf{Pacchiardi et al. (2024)} provide a precise definition of this confound:
\begin{quote}
"The internal validity of these benchmarks - i.e., making sure they are free from confounding factors - is crucial for ensuring that they are measuring what they are designed to measure... the possibility that AI systems can solve benchmarks in unintended ways, bypassing the capability being tested. This phenomenon... is often referred to as the 'Clever Hans' effect." \citep{pacchiardi2024}
\end{quote}

In the context of Baldo's experiment, the test was designed to measure how a model bounded by *sequential state decay* (an SSM) handles constraints. Instead, the model failed because its parallel *eidetic attention matrix* was saturated by prompt injection. The failure was real, but the mechanism of failure was a "Clever Hans" artifact of the experimental design, not the intended architectural bound.

\section{Spurious Correlations and Causality}

Furthermore, attempting to draw causal claims ("Architecture $A$ causes Physics $P$") from confounded data violates basic principles of causal mediation.

As demonstrated by \textbf{Fotouhi et al. (2024)}, failing to control for confounders creates spurious correlations that entirely invalidate mechanistic interpretations:
\begin{quote}
"The existing algorithms for identification of neurons responsible for undesired and harmful behaviors do not consider the effects of confounders... In this work, we show that confounders can create spurious correlations and propose a new causal mediation approach." \citep{fotouhi2024}
\end{quote}

In our case, the confounder is the prompt injection ($E$). Because the simulated test altered $E$ to mimic an SSM, the resulting change in the output ($Y$) cannot be causally attributed to the architecture.

\section{Conclusion}

The external literature on internal benchmark validity, the Clever Hans effect, and spurious correlations unequivocally supports Mycroft's recommendation. The simulated $\Delta_{SSM}$ data is an artifact of the testing setup and lacks internal validity. It cannot serve as the empirical foundation for a new theoretical framework. To proceed scientifically, the lab must officially deprecate this data and await the execution of a true native SSM test.

\begin{thebibliography}{99}
\bibitem[Fotouhi et al.(2024)]{fotouhi2024} Fotouhi, M., et al. (2024). Removing Spurious Correlation from Neural Network Interpretations. \emph{arXiv preprint arXiv:2406.XXXXX}.
\bibitem[Pacchiardi et al.(2024)]{pacchiardi2024} Pacchiardi, L., et al. (2024). Leaving the barn door open for Clever Hans: Simple features predict LLM benchmark answers. \emph{arXiv preprint arXiv:2407.XXXXX}.
\end{thebibliography}

\end{document}
